\setlength{\baselineskip}{20pt}
\chapter{引言}
\label{cha:intro}

\section{研究背景与意义}
视觉判别任务是计算机视觉研究与产业应用的核心方向，广泛应用于检索、安防、交通与智能交互场景。当前主流模型在标准基准上已取得较高精度，但在真实环境中仍受到遮挡、光照变化、天气扰动和域偏移等因素影响，导致特征判别性下降与性能退化。如何在保证部署效率的前提下提升复杂场景鲁棒性，是当前研究中的关键问题。

扩散模型在生成任务中展现出优异的噪声建模能力，但其在判别任务中的潜力尚未被充分挖掘。将扩散去噪机制迁移到判别式表征学习，有望在不依赖额外标注成本的条件下改善特征质量，为分类、检测、分割与检索任务提供统一增强能力。

\section{研究现状与挑战}
现有相关研究主要集中在三条路径：基于数据增强的鲁棒学习、基于结构设计的特征改进、基于扩散思想的任务特定方法。上述方法分别在抗扰动能力、表示能力和建模范式上提供了启发，但仍存在两类共性不足：
\begin{enumerate}
  \item 方法往往面向特定任务或特定数据结构，跨任务复用能力有限；
  \item 若显式引入多步去噪过程，推理时延容易增加，与实际部署需求冲突。
\end{enumerate}

据此，本文聚焦“通用表征增强”和“计算效率保持”两个目标，探索可部署的扩散去噪判别学习方案。

\section{研究内容与主要贡献}
围绕上述问题，本文完成了以下工作：
\begin{enumerate}
  \item 提出基础方法 DenoiseRep，将级联嵌入层与逐层去噪过程统一建模，实现特征提取与去噪协同优化；
  \item 提出参数融合机制，将训练期去噪分支折叠到推理期主干映射中，实现零额外推理开销；
  \item 提出增强方法 DenoiseRep++，引入教师--学生蒸馏与多步融合策略，提升轻量结构下的去噪强度与泛化能力；
  \item 在 ReID、分类、检测、分割和车辆重识别任务上完成系统验证，并将实验结果分别纳入对应方法章节。
\end{enumerate}

\section{论文结构安排}
全文结构如下：第\ref{cha:intro}章为引言；第\ref{cha:chap2}章综述相关工作；第\ref{cha:chap3}章围绕面向无额外推理成本的特征增强算法 DenoiseRep 展开，并在本章内给出实验结果与分析；第\ref{cha:chap4}章在此基础上讨论基于模型蒸馏的增强方法 DenoiseRep++，包含研究点过渡、蒸馏与多步融合策略及实验结果；第\ref{cha:chap5}章总结全文并展望未来工作。
